\documentclass[12pt,a4paper]{article}
\usepackage[margin=2.2cm]{geometry}
\usepackage{amsmath, amssymb, mathtools}
\usepackage{booktabs, array}
\usepackage{hyperref}
\usepackage{xcolor}
\usepackage{microtype}

\hypersetup{
  colorlinks=true,
  linkcolor=blue!50!black,
  urlcolor=blue!50!black,
  citecolor=blue!50!black
}


\title{Explaining the RNN Screenshot\\
\large Why Hidden State Is Memory, and Why the Lower Formula Is Not}
\author{}
\date{\today}

\begin{document}

\maketitle
\tableofcontents

\section{Why the Screenshot Introduces Hidden State}

The screenshot is trying to solve a practical problem, not invent a mysterious symbol. In sequence tasks, the current token rarely makes sense by itself. The word at time \(t\) may depend on something said several steps earlier. If a model stores the entire past literally, the storage burden and the parameter burden both grow quickly.

That is why the screenshot writes
\[
P(x_t \mid x_{t-1}, \ldots, x_1) \approx P(x_t \mid h_{t-1}).
\]
The point is simple: replace a long history by a compact summary \(h_{t-1}\). The hidden state is therefore a \emph{working memory}. It is not the full past. It is a compressed note about the past that is good enough for the next decision.

\section{How the Recurrent Part Works}

The recurrent rule
\[
h_t = f(x_t, h_{t-1})
\]
means that each step uses two inputs:
\begin{itemize}
  \item the new observation \(x_t\),
  \item the previous memory \(h_{t-1}\).
\end{itemize}

The output is a new memory \(h_t\). This is the essential difference from an ordinary feedforward layer. The model does not restart from zero at every step. It carries forward a state.

One may think of \(h_t\) as a notebook that is rewritten one line at a time. The new token writes new information into the notebook, but the old page still influences what remains.

\section{The Lower Formula, Slowly Explained}

The lower half of the screenshot writes
\[
H = \phi(XW_{xh} + b_h), \qquad O = HW_{hq} + b_q.
\]
This part is not yet a recurrent model. It is a standard feedforward map. To see that clearly, it helps to unpack every symbol.

\subsection{What Each Symbol Means}

Assume:
\begin{itemize}
  \item \(X \in \mathbb{R}^{n \times d}\): a batch of inputs,
  \item \(n\): batch size,
  \item \(d\): number of input features per example,
  \item \(W_{xh} \in \mathbb{R}^{d \times h}\): weights from input to hidden features,
  \item \(b_h \in \mathbb{R}^{1 \times h}\): hidden-layer bias,
  \item \(H \in \mathbb{R}^{n \times h}\): hidden representation,
  \item \(W_{hq} \in \mathbb{R}^{h \times q}\): weights from hidden features to output,
  \item \(b_q \in \mathbb{R}^{1 \times q}\): output-layer bias,
  \item \(O \in \mathbb{R}^{n \times q}\): output scores.
\end{itemize}

So the lower formula says nothing more exotic than:
\begin{quote}
turn each input row into hidden features, then turn those hidden features into outputs.
\end{quote}

\subsection{Step 1: From Input Matrix to Hidden Features}

The expression
\[
H = \phi(XW_{xh} + b_h)
\]
contains three small operations.

\paragraph{First, linear mixing.}
Multiply \(X\) by \(W_{xh}\). Each input row with \(d\) numbers becomes a row with \(h\) numbers. This is just feature recombination.

\paragraph{Second, add bias.}
Add \(b_h\). The same bias row is added to every sample in the batch. Bias lets each hidden unit shift its default response.

\paragraph{Third, apply nonlinearity.}
Apply \(\phi\) elementwise, such as ReLU or \(\tanh\). This makes the layer more expressive than a single linear map.

After these three steps, \(H\) is a hidden representation. It is called ``hidden'' only because it is an internal feature layer. It is \emph{not} a time-carrying memory state.

\subsection{Step 2: From Hidden Features to Output}

The second formula
\[
O = HW_{hq} + b_q
\]
maps the hidden features to output scores.

If this is classification, each row of \(O\) is often a vector of logits, and one may later apply softmax. If this is regression, \(O\) may already be the final prediction. In either case, this is still an ordinary one-shot mapping:
\[
\text{input} \rightarrow \text{hidden features} \rightarrow \text{output}.
\]

\subsection{Why This Formula Has No Memory}

The crucial fact is what the formula does \emph{not} contain. It has no term such as \(h_{t-1}\). That means:
\begin{itemize}
  \item it does not explicitly carry information from the previous time step,
  \item it does not update a running state,
  \item it treats each input row as a fresh computation.
\end{itemize}

So even though the letter \(H\) appears, this is still not sequence memory in the RNN sense. It is just a hidden layer in a feedforward network.

This is why the screenshot contrasts the two halves:
\begin{itemize}
  \item the upper half says, ``keep a memory state and update it'';
  \item the lower half says, ``apply a static feature map to the current input''.
\end{itemize}

\section{A Concrete Numerical Example}

Take one input row
\[
x = [2,\;1],
\]
and suppose
\[
W_{xh} =
\begin{bmatrix}
1 & -1\\
2 & 1
\end{bmatrix},
\qquad
b_h = [0,\;1].
\]
Choose \(\phi\) to be ReLU.

Then
\[
xW_{xh}
=
[2,1]
\begin{bmatrix}
1 & -1\\
2 & 1
\end{bmatrix}
=
[4,-1].
\]
After adding bias,
\[
[4,-1] + [0,1] = [4,0].
\]
After ReLU,
\[
H = [4,0].
\]

Now let
\[
W_{hq} =
\begin{bmatrix}
1\\
-1
\end{bmatrix},
\qquad
b_q = [0].
\]
Then
\[
O = HW_{hq} + b_q = [4,0]
\begin{bmatrix}
1\\
-1
\end{bmatrix}
= 4.
\]

This example shows the lower formula very plainly. It only transforms the current input into an output. If we want memory of an earlier word such as \texttt{not}, the formula itself has no dedicated place to store it. That is exactly why the upper recurrent formula introduces \(h_{t-1}\).

\section{Summary}

The screenshot is contrasting two ideas:
\begin{itemize}
  \item RNN: keep a rolling memory \(h_t\) so the past can influence the present,
  \item feedforward network: compute hidden features \(H\) and outputs \(O\) from the current input only.
\end{itemize}

So the lower formula is not difficult because of the symbols. It is only the standard pipeline
\[
\text{input} \rightarrow \text{hidden layer} \rightarrow \text{output layer}.
\]
Its limitation is also standard: by itself, it does not remember earlier time steps.

\end{document}
